\documentclass[12pt, a4paper]{article}

% ------------------------------------------------------------------
% Packages
% ------------------------------------------------------------------
\usepackage[T1]{fontenc}
\usepackage[utf8]{inputenc}
\usepackage{amsmath, amssymb, amsthm}
\usepackage{mathtools}
\usepackage{booktabs}
\usepackage{array}
\usepackage{xcolor}
\usepackage{hyperref}
\usepackage{geometry}
\usepackage{enumitem}
\usepackage{tcolorbox}

\geometry{margin=2.5cm}
\hypersetup{colorlinks=true, linkcolor=blue!60!black, citecolor=blue!60!black}

% ------------------------------------------------------------------
% Theorem environments
% ------------------------------------------------------------------
\newtheoremstyle{condstyle}{}{}{\itshape}{}{\bfseries}{.}{.5em}{}
\theoremstyle{condstyle}
\newtheorem{condition}{Condition}

\theoremstyle{plain}
\newtheorem{proposition}{Proposition}
\newtheorem{corollary}{Corollary}

% ------------------------------------------------------------------
% Coloured boxes
% ------------------------------------------------------------------
\tcbuselibrary{skins}
\newtcolorbox{keyresult}[1][]{enhanced, colback=blue!4!white,
  colframe=blue!40!black, fonttitle=\bfseries, title=#1,
  left=6pt, right=6pt, top=4pt, bottom=4pt, arc=3pt}

\newtcolorbox{notebox}[1][]{enhanced, colback=orange!5!white,
  colframe=orange!60!black, fonttitle=\bfseries, title=#1,
  left=6pt, right=6pt, top=4pt, bottom=4pt, arc=3pt}

\newtcolorbox{warnbox}[1][]{enhanced, colback=red!4!white,
  colframe=red!50!black, fonttitle=\bfseries, title=#1,
  left=6pt, right=6pt, top=4pt, bottom=4pt, arc=3pt}

% ------------------------------------------------------------------
% Notation macros
% ------------------------------------------------------------------
\newcommand{\EN}{\mathbb{E}_N}
\newcommand{\bone}{\mathbf{1}_N}
\newcommand{\bzero}{\mathbf{0}_N}
\newcommand{\IN}{\mathbf{I}_N}
\newcommand{\bD}{\mathbf{D}}
\newcommand{\bY}{\mathbf{Y}}
\newcommand{\bZ}{\mathbf{Z}}
% nuisance vectors
\newcommand{\Yhz}[1]{\hat{\mathbf{Y}}^{z_{#1}}}   % instrument-group outcome pred vector
\newcommand{\Dhz}[1]{\hat{\mathbf{D}}^{z_{#1}}}   % instrument-group treatment pred vector
\newcommand{\Sz}[1]{\mathbf{S}^{z}_{#1}}           % instrument-group smoother matrix
% IPW
\newcommand{\lipwz}[1]{\boldsymbol{\lambda}^{\mathrm{ipw},z}_{#1}}
% pseudo variables
\newcommand{\tY}{\tilde{\mathbf{Y}}}
\newcommand{\tD}{\tilde{\mathbf{D}}}
\newcommand{\tZ}{\tilde{\mathbf{Z}}}
% transformation matrix
\newcommand{\Tiw}{\mathbf{T}^{\mathrm{iv}}}
% estimator
\newcommand{\tiv}{\hat{\tau}^{\mathrm{iv\text{-}aipw}}}
% hat p
\newcommand{\phat}{\hat{p}}

% ------------------------------------------------------------------
\begin{document}
% ------------------------------------------------------------------

\begin{titlepage}
\centering\vspace*{2cm}
{\LARGE\bfseries Wald-AIPW Outcome Weights}\\[0.5em]
{\large A Detailed Derivation in the PIVE Framework of Knaus (2024)}\\[2cm]
{\normalsize Working Notes --- Based on Knaus (2024), Sections 3.2.3 and 4.2.4}
\vfill
\begin{abstract}
\noindent
These notes work through the Wald-AIPW estimator completely:
what it is, why it is built the way it is, how it maps onto the
Pseudo-IV Estimator (PIVE) framework of Knaus (2024), how the
outcome weights are derived from the transformation matrix, and
which of Knaus's six conditions are required for the weights to
exist and to be normalized.
The main result is that standard implementations are only
\emph{scale-normalized} (weights sum to zero) but not
\emph{fully normalized} (treated weights do not sum to one),
because Condition~5b is not fulfilled in practice.
Every algebraic step is shown explicitly, including the
multiplications by $\bone$ and $\bD$ that Knaus (2024) states
only briefly.
\end{abstract}
\end{titlepage}

\tableofcontents\newpage

% ==================================================================
\section{Background: What is Wald-AIPW and why does it exist?}
\label{sec:background}
% ==================================================================

\subsection{The classical Wald estimator in IV settings}

In a binary instrument setting we observe
$(D_i, Z_i, Y_i, X_i)$ where $D_i\in\{0,1\}$ is a binary treatment,
$Z_i\in\{0,1\}$ a binary instrument, $Y_i$ the outcome, and $X_i$ covariates.
The classical Wald (1940) estimator of the Local Average Treatment Effect (LATE) is
\begin{equation}
  \hat\tau^{\mathrm{Wald}}
  = \frac{\EN[Y_i\mid Z_i=1]-\EN[Y_i\mid Z_i=0]}%
         {\EN[D_i\mid Z_i=1]-\EN[D_i\mid Z_i=0]}
  = \frac{\text{reduced form}}{\text{first stage}}.
  \label{eq:wald}
\end{equation}
This is consistent when $Z$ is a valid instrument but ignores covariates entirely.

\subsection{Why augment? The role of AIPW}

Augmented inverse probability weighting (AIPW), introduced by
Robins, Rotnitzky \& Zhao (1994, 1995), improves on simple IPW by
adding a \emph{regression adjustment} term that corrects for
finite-sample imbalance and achieves semiparametric efficiency.
In an unconfoundedness setting (no instrument), the AIPW score for unit $i$ is
\begin{equation}
  \tilde{Y}_i^{\mathrm{aipw}}
  = \hat{Y}^d_{1,i} - \hat{Y}^d_{0,i}
    + \lambda^{\mathrm{ipw}}_{1,i}(Y_i - \hat{Y}^d_{1,i})
    - \lambda^{\mathrm{ipw}}_{0,i}(Y_i - \hat{Y}^d_{0,i}),
  \label{eq:aipw_score}
\end{equation}
where $\hat{Y}^d_{d,i}=\widehat{\mathbb{E}}[Y_i\mid D_i=d, X_i]$ are
treatment-group-specific outcome predictions and
$\lambda^{\mathrm{ipw}}_{1,i}=D_i/\hat{p}(X_i)$,
$\lambda^{\mathrm{ipw}}_{0,i}=(1-D_i)/(1-\hat{p}(X_i))$ are inverse propensity weights.
The score in \eqref{eq:aipw_score} combines the outcome model (first two terms) with
an IPW correction (last two terms), yielding \emph{double robustness}:
the estimator is consistent if either the outcome model or the propensity score model is
correctly specified.

\subsection{Wald-AIPW: applying AIPW to both legs of the Wald ratio}

Tan (2006), building on the Wald structure in \eqref{eq:wald}, applies AIPW
separately to the reduced form ($Y$ on $Z$) and to the first stage ($D$ on $Z$).
Instead of a treatment propensity score, we now use the \emph{instrument} propensity
score $p(X_i)=\mathbb{E}[Z_i\mid X_i]$, estimated as $\phat_i$.
The instrument-group outcome predictions are
\begin{equation}
  \hat{Y}^z_{z,i} = \widehat{\mathbb{E}}[Y_i\mid Z_i=z, X_i], \quad z\in\{0,1\},
  \label{eq:Yhat_z}
\end{equation}
and analogously $\hat{D}^z_{z,i}=\widehat{\mathbb{E}}[D_i\mid Z_i=z,X_i]$.
The instrument-group inverse probability weights are
\begin{equation}
  \lambda^{\mathrm{ipw},z}_{1,i} = \frac{Z_i}{\phat_i},
  \qquad
  \lambda^{\mathrm{ipw},z}_{0,i} = \frac{1-Z_i}{1-\phat_i}.
  \label{eq:ipw_z}
\end{equation}

The Wald-AIPW \emph{reduced form} and \emph{first stage} AIPW scores are:
\begin{align}
  \tilde{Y}^{\mathrm{iv\text{-}aipw}}_i
  &= \hat{Y}^z_{1,i} - \hat{Y}^z_{0,i}
     + \lambda^{\mathrm{ipw},z}_{1,i}(Y_i - \hat{Y}^z_{1,i})
     - \lambda^{\mathrm{ipw},z}_{0,i}(Y_i - \hat{Y}^z_{0,i}),
  \label{eq:Ytilde_iv}\\[4pt]
  \tilde{D}^{\mathrm{iv\text{-}aipw}}_i
  &= \hat{D}^z_{1,i} - \hat{D}^z_{0,i}
     + \lambda^{\mathrm{ipw},z}_{1,i}(D_i - \hat{D}^z_{1,i})
     - \lambda^{\mathrm{ipw},z}_{0,i}(D_i - \hat{D}^z_{0,i}).
  \label{eq:Dtilde_iv}
\end{align}

\begin{notebox}[Key idea]
Compare \eqref{eq:Ytilde_iv} with the plain AIPW score \eqref{eq:aipw_score}:
the structure is \emph{identical}, but treatment group ($d\in\{0,1\}$) is
replaced everywhere by instrument group ($z\in\{0,1\}$), and
treatment IPW weights are replaced by instrument IPW weights.
The two ``legs'' of the Wald ratio in \eqref{eq:wald} are each
augmented the same way that AIPW augments a simple difference in means.
\end{notebox}

The estimator then takes the Wald ratio of these two AIPW scores:
\begin{equation}
  \tiv
  = \frac{\EN[\tilde{Y}^{\mathrm{iv\text{-}aipw}}_i]}%
         {\EN[\tilde{D}^{\mathrm{iv\text{-}aipw}}_i]}.
  \label{eq:wald_aipw}
\end{equation}

% ==================================================================
\section{The PIVE Framework}
\label{sec:pive}
% ==================================================================

Knaus (2024, Definition 1) defines the class of \emph{Pseudo-IV Estimators} (PIVEs)
as those solving an empirical moment condition
\begin{equation}
  \EN\!\left[(\tilde{Y}_i - \hat\tau\,\tilde{D}_i)\,\tilde{Z}_i\right] = 0,
  \label{eq:moment}
\end{equation}
with scalar pseudo-outcome $\tilde{Y}_i$, pseudo-treatment $\tilde{D}_i$, and
pseudo-instrument $\tilde{Z}_i$.
Solving yields
\begin{equation}
  \hat\tau
  = \frac{\EN[\tilde{Z}_i\tilde{Y}_i]}{\EN[\tilde{Z}_i\tilde{D}_i]}
  = (\tilde{\mathbf{Z}}'\tilde{\mathbf{D}})^{-1}\tilde{\mathbf{Z}}'\tilde{\mathbf{Y}}.
  \label{eq:pive_sol}
\end{equation}

\begin{proposition}[Outcome weights of PIVEs, Knaus 2024 Prop.~1]
\label{prop:weights}
If $\tilde{\mathbf{Z}}'\tilde{\mathbf{D}}\neq 0$ and there exists a unique
$N\times N$ transformation matrix $\mathbf{T}$ such that $\mathbf{T}\bY=\tilde{\mathbf{Y}}$,
then
\begin{equation}
  \boldsymbol{\omega}' = (\tilde{\mathbf{Z}}'\tilde{\mathbf{D}})^{-1}\tilde{\mathbf{Z}}'\mathbf{T},
  \quad\text{so that}\quad
  \hat\tau = \boldsymbol{\omega}'\bY = \sum_{i=1}^N \omega_i Y_i.
  \label{eq:prop1}
\end{equation}
\end{proposition}

\noindent
The strategy to derive outcome weights is therefore:
\begin{enumerate}[label=\textbf{Step \arabic*.}]
  \item Write the estimator as a PIVE: identify $\tilde{Y}_i$, $\tilde{D}_i$, $\tilde{Z}_i$.
  \item Find $\mathbf{T}$ such that $\mathbf{T}\bY = \tilde{\mathbf{Y}}$.
  \item Apply Proposition~\ref{prop:weights}.
\end{enumerate}

\medskip\noindent
\textbf{Normalization shortcuts.}
To classify whether outcome weights are fully normalized, Knaus (2024, Section~4.1)
provides three matrix shortcuts that are \emph{sufficient} (not necessary):
\begin{enumerate}[label=\textbf{SC\arabic*.}]
  \item $\mathbf{T}\bone = \bzero$ implies $\sum_i\omega_i = 0$ (weights sum to zero).
  \item $\mathbf{T}\bD = \tilde{\mathbf{D}}$ implies $\sum_i\omega_iD_i = 1$ (treated weights sum to one).
  \item $\mathbf{T}(\bone-\bD)=-\tilde{\mathbf{D}}$ implies $\sum_i\omega_i(1-D_i)=-1$.
\end{enumerate}
When SC1 holds, SC3 follows from SC2 by subtraction.

% ==================================================================
\section{Conditions from Knaus (2024)}
\label{sec:conditions}
% ==================================================================

These are the six conditions Knaus (2024) introduces in Sections~3.1.2 and 4.2.1.
We state each one precisely and discuss its role for Wald-AIPW specifically.

\begin{condition}[Smoother matrix --- existence of outcome weights]\label{cond:C1}
A unique smoother matrix exists producing outcome nuisance predictions as
$\mathbf{S}\bY=\hat{\mathbf{Y}}$. For instrument-group predictions (relevant here):
\begin{equation}
  \textbf{(C1c)}\quad \Sz{1}\bY = \Yhz{1} \quad\text{and}\quad \Sz{0}\bY = \Yhz{0},
  \label{eq:C1c}
\end{equation}
where $(\Sz{z})_{ij} = s^z_{i\leftarrow j}$ is the weight unit $j$'s outcome
receives in predicting unit $i$'s outcome within instrument group $z$.\\[2pt]
\textit{Role:} This is the \textbf{existence condition}. Without a smoother matrix,
$\mathbf{T}$ does not exist and $\boldsymbol{\omega}$ has no closed form.
Methods with non-linear link functions (e.g.\ logistic regression for outcome,
Lasso) typically do not admit a smoother representation.\\[2pt]
\textit{For Wald-AIPW:} C1c must hold; it is assumed but not guaranteed.
In practice, random forests and ridge regression satisfy C1c; Lasso does not.
\end{condition}

\begin{condition}[Covariate matrix with constant]\label{cond:C2}
The covariate matrix $X$ contains a column of ones.\\[2pt]
\textit{Role:} Ensures OLS/TSLS are normalized. Not relevant for Wald-AIPW.
\end{condition}

\begin{condition}[Affine smoother matrix]\label{cond:C3}
All rows of the smoother matrices sum to one:
\begin{equation}
  \Sz{1}\bone = \Sz{0}\bone = \bone.
  \label{eq:C3}
\end{equation}
\textit{Role:} This guarantees $\mathbf{T}\bone = \bzero$ (shortcut SC1), i.e.\ normalized weights.
Most regression smoothers satisfy this: kernel regression, $k$-NN, random forests,
ridge regression, splines. Boosted trees may not.\\[2pt]
\textit{For Wald-AIPW:} C3 is \textbf{the condition for normalization} ($\sum_i\omega_i=0$).
It is met whenever a standard random forest or ridge is used for the outcome models.
\end{condition}

\begin{condition}[No smoothing between treatment groups]\label{cond:C4}
Treatment-group-specific predictions use only observations from that group:
$\mathbf{S}^d_1(\bone-\bD)=\bzero$ and $\mathbf{S}^d_0\bD=\bzero$.\\[2pt]
\textit{Role:} Ensures AIPW is fully normalized. Not directly relevant for Wald-AIPW
(which conditions on instrument groups $z$, not treatment groups $d$).
\end{condition}

\begin{condition}[Outcome smoother applied to treatment]\label{cond:C5}
(C5b, the relevant variant for Wald-AIPW): the instrument-group treatment predictions
are formed using the \emph{same} smoother matrices as the outcome predictions:
\begin{equation}
  \textbf{(C5b)}\quad \Sz{1}\bD = \Dhz{1} \quad\text{and}\quad \Sz{0}\bD = \Dhz{0}.
  \label{eq:C5b}
\end{equation}
\textit{Role:} This guarantees $\mathbf{T}\bD = \tilde{\mathbf{D}}$ (shortcut SC2),
i.e.\ treated weights sum to one.\\[2pt]
\textit{For Wald-AIPW:} C5b is the \textbf{condition for full normalization}.
It requires using the \emph{identical} flexible model (same smoother matrix,
same hyperparameters) for both the outcome nuisance $\hat{Y}^z_{z,i}$ and the
treatment nuisance $\hat{D}^z_{z,i}$.
This is essentially never done in practice --- standard implementations
run separate models for each nuisance, violating C5b.
\end{condition}

\begin{condition}[Normalized IPW weights]\label{cond:C6}
(C6b, the relevant variant): instrument-group IPW weights are Hájek-normalized:
$\lambda^{\mathrm{norm},z}_{z,i} = \lambda^{\mathrm{ipw},z}_{z,i}/\EN[\lambda^{\mathrm{ipw},z}_z]$,
ensuring $\bone'\boldsymbol{\lambda}^{\mathrm{norm},z}_1 = \bone'\boldsymbol{\lambda}^{\mathrm{norm},z}_0 = N$.\\[2pt]
\textit{Role:} Ensures Wald-IPW is fully normalized. For Wald-AIPW, C6b alone
is not sufficient because the regression adjustment terms also matter.
\end{condition}

\begin{notebox}[Summary: which conditions matter for Wald-AIPW]
\begin{tabular}{@{}lll@{}}
\toprule
Condition & Required for & Met in practice? \\
\midrule
C1c & Closed-form weights to exist & Depends on ML method \\
C3  & Weights sum to zero (normalized) & Usually yes (random forest, ridge) \\
C5b & Treated weights sum to one (fully normalized) & Usually \textbf{no} \\
C2, C4, C6 & Not relevant for Wald-AIPW & --- \\
\bottomrule
\end{tabular}
\end{notebox}

% ==================================================================
\section{Step 1: Writing Wald-AIPW as a PIVE}
\label{sec:step1}
% ==================================================================

\subsection{Scalar moment condition}

Following Chernozhukov et al.\ (2018) and Knaus (2024, eq.~18),
the Wald-AIPW moment condition in PIVE form is:
\begin{equation}
  \EN\!\left[
    \left(
      \underbrace{%
        \hat{Y}^z_{1,i}-\hat{Y}^z_{0,i}
        +\lambda^{\mathrm{ipw},z}_{1,i}(Y_i-\hat{Y}^z_{1,i})
        -\lambda^{\mathrm{ipw},z}_{0,i}(Y_i-\hat{Y}^z_{0,i})
      }_{=:\,\tilde{Y}^{\mathrm{iv}}_i}
      -\tiv\cdot
      \underbrace{%
        \hat{D}^z_{1,i}-\hat{D}^z_{0,i}
        +\lambda^{\mathrm{ipw},z}_{1,i}(D_i-\hat{D}^z_{1,i})
        -\lambda^{\mathrm{ipw},z}_{0,i}(D_i-\hat{D}^z_{0,i})
      }_{=:\,\tilde{D}^{\mathrm{iv}}_i}
    \right)\cdot
    \underbrace{1}_{=:\,\tilde{Z}^{\mathrm{iv}}_i}
  \right] = 0.
  \label{eq:moment_iv}
\end{equation}

The PIVE ingredients are therefore:
\begin{align*}
  \tilde{Z}^{\mathrm{iv}}_i &= 1 \quad\Rightarrow\quad \tilde{\mathbf{Z}}' = N^{-1}\bone',\\[3pt]
  \tilde{D}^{\mathrm{iv}}_i &= \hat{D}^z_{1,i}-\hat{D}^z_{0,i}
    +\lambda^{\mathrm{ipw},z}_{1,i}(D_i-\hat{D}^z_{1,i})
    -\lambda^{\mathrm{ipw},z}_{0,i}(D_i-\hat{D}^z_{0,i}),\\[3pt]
  \tilde{Y}^{\mathrm{iv}}_i &= \hat{Y}^z_{1,i}-\hat{Y}^z_{0,i}
    +\lambda^{\mathrm{ipw},z}_{1,i}(Y_i-\hat{Y}^z_{1,i})
    -\lambda^{\mathrm{ipw},z}_{0,i}(Y_i-\hat{Y}^z_{0,i}).
\end{align*}

\subsection{Matrix form of the estimator}

Solving \eqref{eq:moment_iv} and writing in matrix notation (Knaus 2024, eq.~19):
\begin{equation}
  \tiv
  = \underbrace{%
      \left(\bone'\tilde{\mathbf{D}}^{\mathrm{iv}}\right)^{-1}
    }_{(\tilde{\mathbf{Z}}'\tilde{\mathbf{D}})^{-1}}
    \cdot
    \bone'\!
    \underbrace{%
      \left[\Yhz{1}-\Yhz{0}
        +\mathrm{diag}(\lipwz{1})(Y-\Yhz{1})
        -\mathrm{diag}(\lipwz{0})(Y-\Yhz{0})\right]
    }_{=\,\tilde{\mathbf{Y}}^{\mathrm{iv}}},
  \label{eq:wald_aipw_matrix}
\end{equation}
where
\begin{equation}
  \tilde{\mathbf{D}}^{\mathrm{iv}}
  = \Dhz{1}-\Dhz{0}
    +\mathrm{diag}(\lipwz{1})(\bD-\Dhz{1})
    -\mathrm{diag}(\lipwz{0})(\bD-\Dhz{0}).
  \label{eq:Dtilde_iv_vec}
\end{equation}

% ==================================================================
\section{Step 2: Finding the Transformation Matrix $\mathbf{T}^{\mathrm{iv}}$}
\label{sec:step2}
% ==================================================================

We need to find $\Tiw$ such that $\Tiw\bY = \tilde{\mathbf{Y}}^{\mathrm{iv}}$.
This requires Condition~\ref{cond:C1} (C1c): the instrument-group outcome predictions
must be produced by smoother matrices:
$\Sz{1}\bY = \Yhz{1}$ and $\Sz{0}\bY = \Yhz{0}$.

Substituting C1c into $\tilde{\mathbf{Y}}^{\mathrm{iv}}$:
\begin{align}
  \tilde{\mathbf{Y}}^{\mathrm{iv}}
  &= \Sz{1}\bY - \Sz{0}\bY
    + \mathrm{diag}(\lipwz{1})(\bY - \Sz{1}\bY)
    - \mathrm{diag}(\lipwz{0})(\bY - \Sz{0}\bY) \notag\\
  &= \Sz{1}\bY - \Sz{0}\bY
    + \mathrm{diag}(\lipwz{1})\bY - \mathrm{diag}(\lipwz{1})\Sz{1}\bY
    - \mathrm{diag}(\lipwz{0})\bY + \mathrm{diag}(\lipwz{0})\Sz{0}\bY \notag\\
  &= \underbrace{%
      \left[
        \Sz{1} - \Sz{0}
        + \mathrm{diag}(\lipwz{1})(\IN - \Sz{1})
        - \mathrm{diag}(\lipwz{0})(\IN - \Sz{0})
      \right]
    }_{=:\,\Tiw}\bY.
  \label{eq:T_iv}
\end{align}

\begin{corollary}[Outcome weights of Wald-AIPW, Knaus 2024, Corollary~3]
\label{cor:weights}
Under Condition~\ref{cond:C1}(C1c), the outcome weights of Wald-AIPW are:
\begin{equation}
  \boldsymbol{\omega}^{\mathrm{iv}\prime}
  = (\bone'\tilde{\mathbf{D}}^{\mathrm{iv}})^{-1}\cdot
    \bone'\,\Tiw,
  \label{eq:weights_iv}
\end{equation}
where $\Tiw$ is as defined in \eqref{eq:T_iv}.
Per observation:
\begin{equation}
  \boxed{
  \omega^{\mathrm{iv}}_i
  = \frac{1}{\bone'\tilde{\mathbf{D}}^{\mathrm{iv}}}
    \left[
      s^z_{i,1} - s^z_{i,0}
      + \lambda^{\mathrm{ipw},z}_{1,i}(1 - s^z_{i,1})
      - \lambda^{\mathrm{ipw},z}_{0,i}(1 - s^z_{i,0})
    \right],
  }
  \label{eq:w_iv_scalar}
\end{equation}
where $s^z_{i,z}$ is the $i$-th row of $\Sz{z}$ summed (the total weight unit $i$
receives from the instrument-group $z$ smoother).
\end{corollary}

\begin{notebox}[Anatomy of the outcome weight]
Each $\omega^{\mathrm{iv}}_i$ has two parts, divided by the first-stage scalar
$\bone'\tilde{\mathbf{D}}^{\mathrm{iv}}$:
\begin{itemize}
  \item \textbf{Regression part:} $s^z_{i,1} - s^z_{i,0}$ --- how much unit $i$'s
  outcome contributes to the predicted outcome under $Z=1$ minus under $Z=0$.
  This is non-zero for all units (smooth prediction uses entire sample).
  \item \textbf{IPW correction:} $\lambda^{\mathrm{ipw},z}_{1,i}(1-s^z_{i,1})
  - \lambda^{\mathrm{ipw},z}_{0,i}(1-s^z_{i,0})$ --- the residual-based augmentation.
  $\lambda^{\mathrm{ipw},z}_{1,i} = Z_i/\phat_i$ is large for $Z_i=1$ units with low
  propensity score; $\lambda^{\mathrm{ipw},z}_{0,i} = (1-Z_i)/(1-\phat_i)$ is large
  for $Z_i=0$ units with high propensity score.
\end{itemize}
\end{notebox}

% ==================================================================
\section{Step 3: Checking Normalization --- Does $\sum_i\omega_i^{\mathrm{iv}}=0$?}
\label{sec:sc1}
% ==================================================================

We use Shortcut SC1: check whether $\Tiw\bone = \bzero$.
Expanding using \eqref{eq:T_iv}:
\begin{align}
  \Tiw\bone
  &= \Sz{1}\bone - \Sz{0}\bone
    + \mathrm{diag}(\lipwz{1})(\IN-\Sz{1})\bone
    - \mathrm{diag}(\lipwz{0})(\IN-\Sz{0})\bone \notag\\
  &= \Sz{1}\bone - \Sz{0}\bone
    + \lipwz{1} - \mathrm{diag}(\lipwz{1})\Sz{1}\bone
    - \lipwz{0} + \mathrm{diag}(\lipwz{0})\Sz{0}\bone.
  \label{eq:sc1_expand}
\end{align}

\textbf{Apply Condition~\ref{cond:C3} (affine smoother, C3):}
$\Sz{1}\bone = \bone$ and $\Sz{0}\bone = \bone$.
Substituting into \eqref{eq:sc1_expand}:
\begin{align}
  \Tiw\bone
  &\stackrel{\mathrm{C3}}{=}
    \bone - \bone + \lipwz{1} - \mathrm{diag}(\lipwz{1})\bone
    - \lipwz{0} + \mathrm{diag}(\lipwz{0})\bone \notag\\
  &= \lipwz{1} - \lipwz{1} - \lipwz{0} + \lipwz{0} \notag\\
  &= \bzero. \qquad\checkmark
  \label{eq:sc1_result}
\end{align}

\begin{keyresult}[$\tiv$ is normalized under C3]
$\Tiw\bone = \bzero$ under Condition~\ref{cond:C3}, which implies
\[
  \sum_{i=1}^N \omega^{\mathrm{iv}}_i = 0.
\]
\textbf{Why C3?}
The key cancellation in \eqref{eq:sc1_result} is
$\Sz{z}\bone = \bone$, i.e.\ the smoother is affine.
Without C3, $\Sz{1}\bone \neq \bone$ and $\Sz{0}\bone \neq \bone$,
so the terms $(\Sz{1}-\Sz{0})\bone \neq \bzero$ and the cancellation fails.
The IPW part ($\lipwz{1}-\lipwz{0}$) self-cancels \emph{regardless} of C3;
it is the regression part that requires affineness.
This is the same mechanism that normalizes AIPW (Knaus 2024, Section~4.2.3),
directly adapted to instrument groups by swapping superscripts.
\end{keyresult}

% ==================================================================
\section{Step 4: Checking Full Normalization --- Does $\sum_i\omega_i^{\mathrm{iv}}D_i=1$?}
\label{sec:sc2}
% ==================================================================

We use Shortcut SC2: check whether $\Tiw\bD = \tilde{\mathbf{D}}^{\mathrm{iv}}$.
Expanding $\Tiw\bD$ using \eqref{eq:T_iv}:
\begin{align}
  \Tiw\bD
  &= \Sz{1}\bD - \Sz{0}\bD
    + \mathrm{diag}(\lipwz{1})(\IN-\Sz{1})\bD
    - \mathrm{diag}(\lipwz{0})(\IN-\Sz{0})\bD \notag\\
  &= \Sz{1}\bD - \Sz{0}\bD
    + \mathrm{diag}(\lipwz{1})(\bD - \Sz{1}\bD)
    - \mathrm{diag}(\lipwz{0})(\bD - \Sz{0}\bD).
  \label{eq:sc2_expand}
\end{align}

For SC2 to hold, we need this to equal $\tilde{\mathbf{D}}^{\mathrm{iv}}$
as defined in \eqref{eq:Dtilde_iv_vec}:
\begin{equation}
  \tilde{\mathbf{D}}^{\mathrm{iv}}
  = \Dhz{1}-\Dhz{0}
    +\mathrm{diag}(\lipwz{1})(\bD-\Dhz{1})
    -\mathrm{diag}(\lipwz{0})(\bD-\Dhz{0}).
  \label{eq:Dtilde_again}
\end{equation}

Comparing \eqref{eq:sc2_expand} with \eqref{eq:Dtilde_again} term by term:

\begin{center}
\renewcommand{\arraystretch}{1.5}
\begin{tabular}{@{}lll@{}}
\toprule
Term in $\Tiw\bD$ & Term in $\tilde{\mathbf{D}}^{\mathrm{iv}}$ & Equal when \\
\midrule
$\Sz{1}\bD$ & $\Dhz{1}$ & C5b: $\Sz{1}\bD = \Dhz{1}$ \\
$-\Sz{0}\bD$ & $-\Dhz{0}$ & C5b: $\Sz{0}\bD = \Dhz{0}$ \\
$\mathrm{diag}(\lipwz{1})(\bD - \Sz{1}\bD)$ & $\mathrm{diag}(\lipwz{1})(\bD-\Dhz{1})$ & C5b \\
$-\mathrm{diag}(\lipwz{0})(\bD - \Sz{0}\bD)$ & $-\mathrm{diag}(\lipwz{0})(\bD-\Dhz{0})$ & C5b \\
\bottomrule
\end{tabular}
\end{center}

\medskip
All four terms agree if and only if Condition~\ref{cond:C5}(C5b) holds,
i.e.\ $\Sz{z}\bD = \Dhz{z}$ for $z\in\{0,1\}$.
\textbf{If C5b holds:}
\begin{align}
  \Tiw\bD
  &\stackrel{\mathrm{C5b}}{=}
    \Dhz{1} - \Dhz{0}
    + \mathrm{diag}(\lipwz{1})(\bD-\Dhz{1})
    - \mathrm{diag}(\lipwz{0})(\bD-\Dhz{0})
  = \tilde{\mathbf{D}}^{\mathrm{iv}}. \qquad\checkmark
  \label{eq:sc2_c5b}
\end{align}

This establishes SC2: $\Tiw\bD = \tilde{\mathbf{D}}^{\mathrm{iv}}$, and since we already
showed SC1 holds (under C3), we then have
\[
  \boldsymbol{\omega}^{\mathrm{iv}\prime}\bD
  = (\bone'\tilde{\mathbf{D}}^{\mathrm{iv}})^{-1}\bone'\Tiw\bD
  = (\bone'\tilde{\mathbf{D}}^{\mathrm{iv}})^{-1}\bone'\tilde{\mathbf{D}}^{\mathrm{iv}}
  = 1. \qquad\checkmark
\]

\begin{keyresult}[$\tiv$ is fully normalized under C3 and C5b]
If both Condition~\ref{cond:C3} and Condition~\ref{cond:C5}(C5b) hold, then
\[
  \sum_i\omega^{\mathrm{iv}}_i = 0,\quad
  \sum_i\omega^{\mathrm{iv}}_iD_i = 1,\quad
  \sum_i\omega^{\mathrm{iv}}_i(1-D_i) = -1 \quad\text{(by subtraction)}.
\]
\end{keyresult}

% ==================================================================
\section{Step 5: Why C5b Fails in Practice}
\label{sec:c5b_fail}
% ==================================================================

Condition~\ref{cond:C5}(C5b) requires that the \emph{same} smoother matrix
$\Sz{z}$ used to predict outcomes is also used to predict treatments:
\[
  \Sz{1}\bY = \Yhz{1} \quad\text{and}\quad \Sz{1}\bD = \Dhz{1},
\]
and analogously for $z=0$.
In concrete terms, this means fitting one model to the $Z=1$ subsample
and using its predictions for \emph{both} the outcome and the treatment.

This is at odds with standard ML practice, where a separate model is fit for
each target variable.
For example, in the DoubleML package (Bach et al.\ 2024) or the grf package,
the practitioner specifies:
\begin{itemize}
  \item A learner for $\hat{Y}^z_{1,i}$: e.g.\ random forest with outcome as target.
  \item A learner for $\hat{D}^z_{1,i}$: e.g.\ random forest with treatment as target.
\end{itemize}
These two learners produce \emph{different} smoother matrices
${\Sz{1}}^{Y} \neq {\Sz{1}}^{D}$, since the split points in the trees differ.
Therefore, ${\Sz{1}}^{Y}\bD \neq \hat{\mathbf{D}}^{z_1}$ in general, violating C5b.

\begin{warnbox}[Consequence of C5b failure]
When C5b does not hold, $\Tiw\bD \neq \tilde{\mathbf{D}}^{\mathrm{iv}}$.
This means:
\begin{itemize}
  \item SC2 fails: $\sum_i\omega^{\mathrm{iv}}_iD_i \neq 1$ in general.
  \item The Wald-AIPW outcome weights are only \textbf{scale-normalized}
  (sum to zero, treated and untreated sums equal some constant $c\neq 1$ and $-c$).
  \item The estimate is not guaranteed to be translation invariant:
  adding a constant $c$ to all $Y_i$ does not necessarily cancel out.
\end{itemize}
This is exactly what Knaus (2024) documents empirically in the EMCS
(Figure~1): Wald-AIPW with random forests deviates from the true effect of 1
in a setting where fully-normalized estimators recover it exactly.
\end{warnbox}

\subsection{What C5b would look like concretely}

Suppose we used ridge regression for both outcome and treatment within
each instrument group. The ridge smoother for the $Z=1$ subsample is
$\Sz{1} = \mathbf{X}_1(\mathbf{X}_1'\mathbf{X}_1 + \lambda\mathbf{I})^{-1}\mathbf{X}_1'$,
where $\mathbf{X}_1$ contains covariates of $Z=1$ units only.
Using the same matrix for both:
\[
  \Yhz{1} = \Sz{1}\bY, \qquad \Dhz{1} = \Sz{1}\bD.
\]
This satisfies C5b and makes Wald-AIPW fully normalized.
However, this forces identical functional form and regularization for
outcome and treatment --- an unrealistic constraint when their
distributions differ substantially (e.g.\ continuous outcome vs.\ binary treatment).

% ==================================================================
\section{Step 6: Untreated Weights and Summary}
\label{sec:summary}
% ==================================================================

\subsection{Untreated weights sum to minus one (under C3 and C5b)}

Under C3 and C5b, we have $\sum_i\omega^{\mathrm{iv}}_i = 0$ and
$\sum_i\omega^{\mathrm{iv}}_iD_i = 1$, so by subtraction:
\[
  \sum_i\omega^{\mathrm{iv}}_i(1-D_i)
  = \sum_i\omega^{\mathrm{iv}}_i - \sum_i\omega^{\mathrm{iv}}_iD_i
  = 0 - 1 = -1. \qquad\checkmark
\]

\subsection{Contrast with AIPW}

It is instructive to compare why \emph{AIPW} (treatment-based) achieves full normalization
in standard implementations while \emph{Wald-AIPW} (instrument-based) does not.

For AIPW, the treatment predictions $\hat{D}$ do not enter the moment condition at all ---
the pseudo-treatment is just $\tilde{D}_i^{\mathrm{aipw}} = 1$.
Full normalization of AIPW requires only C3 and C4 (no between-group smoothing):
the estimator is \emph{self-normalizing} because fitting separate outcome models within
each treatment group (C4) automatically produces the right structure.
There is no analogue of C5 to worry about.

For Wald-AIPW, however, the first stage (pseudo-treatment $\tilde{D}^{\mathrm{iv}}_i$)
contains its own nuisance predictions $\hat{D}^z_{z,i}$.
These sit inside the denominator of the weight expression, and for SC2 to hold
they must be generated by the same smoother as $\hat{Y}^z_{z,i}$ in the numerator.
This creates the C5b requirement that AIPW simply does not face.

\subsection{Overall classification table}

\begin{table}[h!]
\centering
\caption{Conditions and normalization properties for Wald-AIPW and related estimators}
\label{tab:summary}
\renewcommand{\arraystretch}{1.4}
\begin{tabular}{@{}lllll@{}}
\toprule
Estimator & Closed-form & $\sum\omega_i=0$ & $\sum\omega_iD_i=1$ & $\sum\omega_i(1-D_i)=-1$ \\
\midrule
Wald-AIPW & C1c & C3 & C3 \& C5b & C3 \& C5b \\
Wald-RA   & C1c & C3 & C3 \& C5b & C3 \& C5b \\
Wald-IPW  & always & C6b & C6b & C6b \\
AIPW      & C1b & C3 & C3 \& C4 & C3 \& C4 \\
\midrule
\multicolumn{5}{@{}l@{}}{\small\textit{In practice:} C3 usually holds; C5b does \textbf{not}; hence standard Wald-AIPW is only scale-normalized.}\\
\bottomrule
\end{tabular}
\end{table}

% ==================================================================
\section{Why Does This Matter? Connection to the EMCS}
\label{sec:emcs}
% ==================================================================

Knaus (2024, Section~4.4) runs an Empirical Monte Carlo Study with a synthetic
outcome $Y^*_i = 1 + D_i$, where the true treatment effect is exactly 1.
For any fully-normalized estimator, $\boldsymbol{\omega}'\mathbf{Y}^* = \boldsymbol{\omega}'(\bone + \bD) = 1$,
since $\boldsymbol{\omega}'\bone = 0$ and $\boldsymbol{\omega}'\bD = 1$.
Non-normalized estimators will deviate from 1 because:
\begin{align*}
  \boldsymbol{\omega}^{\mathrm{iv}\prime}\mathbf{Y}^*
  &= \boldsymbol{\omega}^{\mathrm{iv}\prime}(\bone + \bD)
   = \underbrace{\boldsymbol{\omega}^{\mathrm{iv}\prime}\bone}_{=0\text{ (C3)}}
   + \underbrace{\boldsymbol{\omega}^{\mathrm{iv}\prime}\bD}_{\neq 1\text{ (C5b fails)}}.
\end{align*}
The deviation from 1 is entirely due to $\boldsymbol{\omega}^{\mathrm{iv}\prime}\bD \neq 1$,
which is the direct algebraic consequence of C5b failing.
When XGBoost (a non-affine smoother) is used, even C3 fails, so
$\boldsymbol{\omega}^{\mathrm{iv}\prime}\bone \neq 0$ as well --- explaining the
extreme outliers (estimates between $-16$ and $55$) seen in the EMCS figure
for Wald-AIPW with XGBoost.

% ==================================================================
\section*{References}
% ==================================================================

\begin{description}
  \item[Bach et al.\ (2024)] DoubleML: An object-oriented implementation of double
  machine learning in R. \textit{Journal of Statistical Software}, 108(3), 1--56.
  \item[Chernozhukov et al.\ (2018)] Double/Debiased machine learning for treatment
  and structural parameters. \textit{The Econometrics Journal}, 21(1), C1--C68.
  \item[Knaus (2024)] Treatment effect estimators as weighted outcomes.
  arXiv:2411.11559v2. University of Tübingen.
  \item[Robins, Rotnitzky \& Zhao (1994, 1995)] Estimation of regression coefficients
  when some regressors are not always observed / Analysis of semiparametric regression
  models for repeated outcomes in the presence of missing data.
  \textit{Journal of the American Statistical Association}, 89/90.
  \item[Tan (2006)] Regression and weighting methods for causal inference using
  instrumental variables. \textit{JASA}, 101(476), 1607--1618.
\end{description}

\end{document}
